Computers and digital display devices have become widely used in learning, office work, and remote working, resulting in prolonged screen exposure among students and employees~\cite{ref1,ref2}. Prolonged computer use has been associated with both visual symptoms and work-related musculoskeletal disorders, with the neck, shoulders, upper back, and lower back among the body regions most frequently affected~\cite{ref1,ref2}. Factors such as extended computer-use duration, static sitting, and inappropriate working posture have also been identified as relevant contributors to musculoskeletal problems among computer users~\cite{ref2}. In addition, deficiencies in workstation arrangement, monitor position, and lighting conditions may coexist with both musculoskeletal symptoms and visual discomfort~\cite{ref3}. These findings suggest that computer-use ergonomics should not be evaluated based on a single factor alone.

Among posture-related factors, head angle (HA) and viewing distance (VD) are two particularly relevant indicators. Hashimoto et al. reported that HA and VD were significantly associated with sagittal spinal alignment in computer users, especially in the cervical C2--C7 region~\cite{ref4}. Their findings also showed that a more forward spinal alignment tended to be accompanied by a larger head angle and a shorter viewing distance~\cite{ref4}. Earlier, Rempel et al. demonstrated that changes in display distance could lead to changes in head and trunk position during computer tasks~\cite{ref5}. Similarly, a study involving 309 Japanese office workers found that an inappropriate eye-to-monitor distance, particularly below 40~cm, was associated with greater neck pain intensity~\cite{ref6}. These findings support the use of head angle and viewing distance as quantitative indicators of postural ergonomics.

In addition to posture, ambient lighting is an important component of visual ergonomics. Visual conditions at computer workstations are influenced by factors such as illuminance, luminance, lighting direction, glare, correlated color temperature, and display characteristics~\cite{ref7}. Experimental evidence has also shown that changes in ambient illumination can affect visual fatigue during screen-based tasks~\cite{ref8}. However, a low-cost IoT system cannot assess all aspects of lighting quality. Therefore, the present study uses ambient light level as a simple environmental indicator that can be measured directly using a light sensor, rather than as a complete measure of lighting quality.

Another potentially useful visual indicator is blink rate. Chidi-Egboka et al. reported that blink rate decreased significantly during reading tasks, including tasks performed on digital devices, compared with non-reading activities~\cite{ref9}. More importantly, Lapa et al. demonstrated that blink rate could be detected and monitored using a conventional computer camera under real-life conditions; in their exploratory study, lower blink rates were associated with higher Computer Vision Syndrome scores~\cite{ref10}. However, blink rate can also vary according to task characteristics and cognitive demand~\cite{ref11}. Therefore, the present study treats blink rate as a supporting visual ergonomic indicator, rather than as an independent diagnostic measure.

Taken together, head angle, viewing distance, ambient light level, and blink rate represent different but complementary aspects of computer-use ergonomics. Head angle and viewing distance mainly reflect postural ergonomics, whereas blink rate and ambient light level provide additional information related to visual ergonomics. In addition, exposure duration should also be considered because an unfavorable posture or condition lasting only a few seconds is not equivalent to the same condition being sustained for a prolonged period~\cite{ref2,ref12}. A multimodal assessment that incorporates temporal information may therefore provide a more complete representation of the user's ergonomic condition than a single-threshold warning.

Several technologies have been developed to automate ergonomic monitoring. Wearable-sensor systems can provide continuous motion-related information, but they require users to attach and maintain additional devices on the body~\cite{ref13}. For example, ERG-AI used data from five tri-axial accelerometers together with machine learning and large language models to predict occupational postures and generate health-related recommendations~\cite{ref14}. Another approach is to integrate sensors into furniture. SitWell used two $32 \times 32$ pressure-sensing mats embedded in a chair to classify 19 sitting postures, monitor posture duration, and generate personalized feedback using GPT-4o~\cite{ref15}. Although these systems demonstrate the potential of combining sensing and AI, they still depend on wearable devices or specialized hardware~\cite{ref14,ref15}.

Computer vision provides a non-contact and potentially lower-cost alternative. RGB cameras combined with pose-estimation models can be used to estimate head, shoulder, and upper-body posture~\cite{ref16}. However, the accuracy of computer-vision-based ergonomic assessment depends on factors such as camera position and the selected pose-estimation model~\cite{ref16}. Cameras can also be used to monitor visual behavior such as blink rate~\cite{ref10}. This creates the possibility of using a single camera for both postural monitoring and visual-behavior monitoring, without requiring separate eye-tracking or wearable devices.

Although the individual components of the problem have been investigated, current studies generally examine posture, viewing distance, lighting, blink behavior, and AI-based feedback in separate research contexts or systems~\cite{ref4,ref7,ref10,ref14,ref15,ref16}. Recent intelligent ergonomic systems have begun to move from simple detection toward personalized recommendations~\cite{ref14,ref15}, but they commonly rely on wearable sensors or smart furniture. Therefore, there remains limited research on a compact, low-cost, non-contact IoT-AI system that simultaneously integrates head angle, viewing distance, blink rate, ambient light level, and exposure duration, and then uses the combined information to generate context-aware recommendations.

This gap is important because the same detected posture may require different responses depending on the surrounding context. For example, a user with a large head angle but an appropriate viewing distance and a short exposure duration may not require the same level of intervention as a user who simultaneously maintains a large head angle, sits too close to the screen, works under unsuitable lighting, and remains in that condition for a prolonged period. Therefore, rather than simply classifying posture as correct or incorrect, a multimodal system may provide a more meaningful interpretation of why the condition is unfavorable and how it should be corrected.

To bridge these limitations, this research introduces an Edge Vision IoT--LLM framework, a cost-effective, multimodal system centered on the \textbf{Freenove ESP32-S3} microcontroller and the \textbf{OV3660} CMOS image sensor. The system leverages \textbf{MediaPipe} and \textbf{OpenCV} algorithms to perform real-time 3D head pose estimation (Pitch, Yaw, and Roll) alongside blink rate detection via Eye Aspect Ratio (EAR) analysis. Simultaneously, an ultrasonic sensor and a Light Dependent Resistor (LDR) provide precise telemetry for viewing distance and ambient illuminance, respectively. By synthesizing these data streams with cumulative exposure duration, the system conceptualizes the user's ergonomic well-being through two complementary dimensions:

\begin{enumerate}
    \item \textbf{Postural Ergonomic State:} Evaluated based on head pose orientation (Pitch, Yaw, Roll), eye-to-screen distance, and postural persistence.
    \item \textbf{Visual Ergonomic State:} Evaluated based on ocular blink frequency (EAR-derived), workstation illuminance levels, and continuous screen-time duration.
\end{enumerate}

The integrated multimodal state is processed by a \textbf{Cloud-based Large Language Model (LLM)} reasoning engine to generate personalized, context-aware corrective recommendations. It is important to note that the proposed framework is designed as a preventive screening and behavioral support tool; it is not intended for the clinical diagnosis of musculoskeletal disorders, Computer Vision Syndrome (CVS), or chronic ocular diseases. The primary objective of this study is to develop and validate a low-latency, non-intrusive IoT--AI system that fuses postural, spatial, and environmental metrics to provide real-time feedback and mitigate digital-related health risks.

\subsection{Research Questions}
\label{subsec:research_questions}

To achieve the primary objective and evaluate the performance of the proposed system, this study addresses the following four research questions:
\begin{itemize}
    \item \textbf{RQ1:} How accurately can the ESP32-CAM-based system estimate head angle/posture and blink rate during computer use?
    \item \textbf{RQ2:} How accurately and consistently can the distance and light sensors measure viewing distance and ambient light level under different experimental conditions?
    \item \textbf{RQ3:} Can the integration of head angle/posture, viewing distance, blink rate, ambient light level, and exposure duration distinguish different visual-postural ergonomic conditions?
    \item \textbf{RQ4:} Can AI-generated recommendations based on multimodal ergonomic data provide relevant, clear, useful, and actionable guidance for correcting unfavorable ergonomic conditions?
\end{itemize}

\subsection{Expected Contributions}
\label{subsec:expected_contributions}

The main contributions of this study are summarized as follows:
\begin{enumerate}
    \item A low-cost, non-contact multimodal ergonomic monitoring architecture based on an ESP32-CAM combined with distance and ambient light sensors.
    \item A comprehensive monitoring approach that integrates postural, visual, environmental, and temporal indicators rather than relying on a single variable.
    \item An evaluation framework for jointly considering postural ergonomics and visual ergonomics in a unified system.
    \item A mechanism for generating context-aware personalized corrective recommendations using AI, moving beyond simple single-threshold warnings.
\end{enumerate}

